\chapter{结论与展望}\label{Chap:chap6}

\section{研究总结}\label{Sec6.1}

本文提出Swift算法，一种面向远距离传输、多接入链路和异构终端应用条件的启发式多信号融合自适应网络拥塞控制方法。Swift通过协议栈外的确定性规则在确认或拥塞事件上计算窗口决策，不依赖训练样本与神经网络推理；在给定仿真配置和随机种子时，其决策过程可按规则分支检查并重复执行。

本文的研究历程可以概括为以下几个阶段：

\textbf{问题分析阶段}：首先深入分析了远距离传输、蜂窝/无线/卫星等多接入链路以及手机、笔记本、台式机和边缘设备等异构终端并存场景中的拥塞控制挑战。在此基础上，系统性地分析了传统拥塞控制算法（基于丢包的和基于延迟的）的优缺点、ECN机制的潜力与利用不足，以及学习型拥塞控制方法在训练依赖、推理开销与可审计性方面的工程障碍，确定了以协议栈内部信号驱动的白盒启发式设计作为研究方向。

\textbf{方案设计阶段}：基于问题分析的结论，设计了Swift算法的核心组件。从状态设计入手，提出了11维有效观测子空间（嵌入于15元素状态传输容器，另含4项元数据通道）；然后设计了多信号融合拥塞检测方法，采用分层优先级机制区分超时、ECN驱动的窗口缩减与普通丢包；接着设计了差异化拥塞响应机制和启发式反馈自适应参数调优机制；最后将各组件整合为以$\alpha \times BDP$为目标的融合控制策略。

\textbf{实现与验证阶段}：在ns-3.40网络仿真平台上实现了Swift算法原型，仿真进程与决策进程经ZeroMQ同步请求—应答信道交互；瓶颈网关配置启用ECN标记的RED主动队列管理，并对四种被比较协议对称开启ECN。实验矩阵覆盖近端高速、共享收敛链路、无线局域网接入、蜂窝移动接入、城域与长途/跨域广域网及LEO/GEO卫星通信共36个场景，在"纯TCP"与"TCP $+$ 平均负载为瓶颈32\%（峰值64\%、50\%占空比）的UDP Burst"两类设置下对四种协议各运行一次，共288次仿真。指标由ns-3 FlowMonitor记录的前向TCP数据流计算，并通过流数完整性、物理取值范围和配置一致性检查。

\textbf{分析阶段}：Swift在多数纯TCP与UDP突发场景中呈现较高聚合吞吐量，相对逐场景最强基线的增益约为2\%$\sim$6\%。时延、丢包与公平性则具有场景差异：纯TCP下Swift相对BBR的平均时延高约0.8\%，个别场景最高约10.4\%；UDP突发下相对最低时延基线的最大差异约12.2\%；丢包率与NewReno相当，公平性总体接近1但并非逐场景最优。据此，本文将排队时延感知的$\alpha$调度、多时间尺度BDP估计、多种子重复和机制消融列为后续工作。

本文的主要贡献和研究成果总结如下：

\textbf{（1）多信号拥塞状态感知}

本文设计了11维有效观测子空间（嵌入于15元素状态传输容器，另含4项跨进程路由用的元数据通道），将显式拥塞通知（ECN）状态、拥塞控制状态机状态、拥塞事件类型等TCP协议栈内部信息与窗口、在途字节和RTT观测结合。该状态集合为超时、ECN驱动的窗口缩减和普通丢包三类语义判定提供输入，但不代表覆盖协议栈的全部状态。

\textbf{（2）启发式反馈自适应的融合控制}

本文按滑动时间窗口内累计确认字节的差分计算交付速率，经容量40的窗口最大值滤波并乘以最小RTT得到BDP估计。控制器依据最小RTT感知的RTT膨胀阈值、性能反馈值快线EMA相对慢线基线EMA的偏移以及连续增长趋势，调节每流乘性增加因子$\alpha\in[0.85,1.30]$。慢启动阶段结合HyStart风格RTT膨胀检测；拥塞避免阶段围绕$\alpha\times BDP$目标窗口以有界步长拉升或回落。该控制过程由显式规则执行，不包含策略训练或神经网络推理。

\textbf{（3）面向不确定拥塞信号的稳定性与安全约束}

Swift对普通丢包、ECN和超时分别采用0.70、0.75和0.50的窗口保留因子。连续缩减计数在超过3次后保持当前窗口，显式降窗后冻结4个ACK事件中的窗口增长；Python控制器将拥塞窗口限制在$[4\times MSS,\max(4\times BDP,200\times MSS)]$内，环境侧另执行$2\times MSS$下界校验。当协议栈进入CA\_LOSS时，尚未应用的窗口决策被作废。上述机制用于限制连续缩减、快速反弹、窗口越界和陈旧动作覆盖，其独立贡献仍需消融实验验证。

在实验验证方面，本文在ns-3.40网络仿真平台上实现了Swift算法，并在36个场景、288次仿真（两种设置各144次）的全矩阵下进行了系统性对比实验。实验结果表明：

\begin{itemize}
    \item 纯TCP设置下，Swift在多数场景呈现较高聚合吞吐量，相对逐场景最强基线的增益范围为$+2.1\%\sim+5.8\%$（均值$+4.1\%$）；该趋势覆盖近端高速、无线/蜂窝接入及LEO/GEO卫星、长途与跨域WAN等路径。
    \item 平均单向时延与基线同量级：相对NewReno/CUBIC均值平均低约2.3\%（36个场景中24个更低），相对BBR平均高约0.8\%，纯TCP个别场景最高约10.4\%；UDP突发下相对最低时延基线的最大差异约12.2\%。
    \item Swift平均丢包率与NewReno相当且不具有逐场景一致优势；Jain指数均值为0.9989（纯TCP）和0.9990（UDP突发），总体接近1并与基线相当，分别仅在11个和17个场景严格最优。
    \item GEO卫星、LEO卫星、长途WAN与跨域WAN等代表性远距离场景呈现吞吐优势；其中GEO场景相对最强基线的增益为4.2\%（纯TCP）和3.3\%（UDP突发）。
    \item 平均负载为瓶颈32\%的UDP突发使四协议吞吐保持率均值落在68.6\%--68.8\%；Swift在多数配对场景呈现$+2.2\%\sim+5.9\%$的绝对吞吐优势，但保持率与基线基本一致。
    \item 当前长RTT结果未再呈现旧版逐ACK投递速率估计器对应的明显低吞吐特征，这与时间窗口估计器的设计目标一致；由于缺少新旧估计器消融，本文不把该观察表述为单机制因果证明。
\end{itemize}

\section{研究展望}\label{Sec6.2}

Swift在当前36场景矩阵中多数呈现吞吐优势，但时延、丢包与公平性结果并不一致，且结论受限于ns-3仿真环境、单一随机种子、同步决策路径与RED/ECN配置范围。以下内容作为后续研究方向，不作为本文已经完成的实现或实验结论。

\subsection{参数自适应与窗口控制的细化}

实验记录了两处可继续优化的边界：（1）Swift在个别深缓冲/低速率场景的相对时延抬升最高约10\%；（2）长RTT路径上的吞吐增益（+2.2\%$\sim$+4.8\%）低于近端场景。可探索的方向包括：

\begin{itemize}
    \item 排队延迟显式估计：引入$q_{est} = RTT_{smoothed} - RTT_{min}$作为$\alpha$调整的额外输入，使$\alpha$在排队延迟持续增长时自动收敛至1.0以下。
    \item 多时间尺度BDP估计：区分传播延迟和排队延迟的贡献，避免将排队延迟计入BDP估计。
    \item 受约束的带宽探测：在不突破当前$\alpha\in[0.85,1.30]$稳定性边界的前提下，为长RTT路径设计更细粒度的窗口拉升规则。
    \item 机制消融实验：分别关闭ECN/丢包差异化保留因子、连续缩减保护与降窗冻结，量化各机制的独立贡献；当前数据无法完成该分离。
\end{itemize}

\subsection{策略表达能力扩展}

当前Swift算法采用基于规则的启发式决策机制，具有可解释性和低延迟优势，但表示能力受人工规则结构限制。后续研究可在不改变本文实验结论的前提下，探索离线强化学习训练神经网络策略、再通过模型蒸馏提取为紧凑规则的路线；也可引入元学习方法加快不同网络环境之间的参数迁移。相关方法需要在完整实验矩阵、异常过滤和可复现日志基础上重新评估，在获得结果之前不作为本文系统的组成部分。

\subsection{多目标优化扩展}

当前Swift的优化目标主要是高吞吐量、低延迟和低丢包。在移动终端、边缘计算和多业务并发场景中，后续研究可进一步引入公平性、能耗、可靠性和业务优先级等目标。多目标优化存在吞吐、延迟、公平性和能耗之间的冲突，需要基于约束优化方法重新设计反馈函数与评价体系。

\subsection{真实网络部署与验证}

本文实验主要在ns-3仿真平台上进行，与真实网络环境存在差距。后续部署验证需要将算法实现为Linux内核TCP模块或QUIC用户态协议栈，并在覆盖近端高速、蜂窝/无线接入、广域网与卫星链路的真实或半实物测试床上重新测量。真实环境中的NIC offload、中断合并、NUMA效应、操作系统调度、无线链路波动和终端算力差异均可能影响实际性能，因此不能直接将本文仿真结果外推为生产部署效果。

\subsection{实验方法的扩展}

受单次确定性运行限制，本文未报告跨运行号置信区间。后续应使用多个ns-3.40 \texttt{RngRun}对关键场景重复仿真并给出均值与最坏值；同时可扩展队列管理策略（丢尾、CoDel）与突发强度/占空比扫描，检验结论对配置的稳健性。

\subsection{安全性、鲁棒性与跨层协同}

后续研究需要进一步验证算法在异常输入、通信中断、状态损坏和恶意跨流量下的行为边界。可行方向包括：输入合法性检查、异常状态回退、保守策略兜底、形式化安全性质描述，以及与应用层、链路层和操作系统调度器的协同机制。上述方向均需要独立实现和实验验证，本文不将其作为当前系统已经具备的能力。

\subsection{异构网络与新兴传输协议适配}

远距离传输和异构终端场景通常跨越有线、无线、蜂窝、卫星和边缘接入等多类链路。后续研究可进一步考察无线非拥塞丢包识别、卫星链路长RTT适配、5G突发抖动处理、终端算力约束下的轻量化执行，以及将多信号融合、差异化响应和自适应调参思想迁移到QUIC用户态协议栈中的可行性。相关扩展应基于新的实现版本、完整实验矩阵和异常过滤流程重新评估。

\section{应用场景展望}\label{Sec6.3}

基于本文仿真结果，Swift更适合作为面向远距离传输、异构终端和多接入链路的拥塞控制研究原型，而不是已经完成真实网络生产部署验证的协议。其潜在应用场景包括：

\begin{itemize}
    \item \textbf{跨地域内容传输}：在长RTT路径上维持更高吞吐并控制队列驻留（跨域/长途WAN与GEO/LEO场景实测领先）。
    \item \textbf{蜂窝与无线接入}：利用ECN、RTT和在途字节等多源信号提高对接入链路波动的适应能力。
    \item \textbf{移动终端与边缘应用}：面向手机、笔记本、桌面终端和边缘设备并存的场景，探索更稳健的参数自适应机制。
    \item \textbf{卫星与广域链路}：在高传播时延和高BDP路径上保持带宽估计有效并减少拥塞恢复代价。
\end{itemize}

上述应用方向体现了算法设计的潜在价值，但不意味着本文仿真结果可以直接等同于生产网络收益。实际部署前仍需要完成内核或用户态协议栈实现、真实测试床验证、跨流公平性评估和长期稳定性评估。

\section{研究意义与社会价值}\label{Sec6.4}

本研究的理论意义主要体现在三个方面：首先，提出了面向远距离和异构终端传输的11维有效观测子空间，为协议栈信号驱动的拥塞控制状态设计提供了参考；其次，提出了多信号融合拥塞检测与差异化响应机制，为区分ECN、丢包和超时等不同拥塞信号提供了可解释规则；最后，提出了结合BDP估计、路径感知$\alpha$调节、降窗后冻结和连续缩减保护的白盒启发式融合控制框架，验证了"确定性规则+轻量反馈自适应"设计点在跨场景一致性上的可行性。

本研究通过ns-3.40构建了一个可重复执行的实验原型，考察多信号感知与启发式反馈自适应在长距离、多接入和异构终端应用条件下的可行性。在RED/ECN对称配置中，该方法在多数纯TCP与UDP突发场景呈现吞吐优势，时延、丢包和公平性总体与基线相当但存在场景差异。由于实验采用链路参数抽象而非真实终端硬件，实际部署效果仍需在真实协议栈和网络环境中验证。
